%% LyX 2.4.4 created this file.  For more info, see https://www.lyx.org/.
%% Do not edit unless you really know what you are doing.
\documentclass[oneside,english]{amsart}
\usepackage[T1]{fontenc}
\usepackage[latin9]{inputenc}
\usepackage{amsthm}
\usepackage{amssymb}
\usepackage{graphicx}
\usepackage{geometry}
\geometry{verbose,tmargin=1in,bmargin=1in,lmargin=1in,rmargin=1in}
\usepackage[authoryear]{natbib}

\makeatletter
%%%%%%%%%%%%%%%%%%%%%%%%%%%%%% Textclass specific LaTeX commands.
\numberwithin{equation}{section}
\numberwithin{figure}{section}

%%%%%%%%%%%%%%%%%%%%%%%%%%%%%% User specified LaTeX commands.
\usepackage{hyperref}
\usepackage[usenames,dvipsnames]{xcolor} %adds additional color options
\hypersetup{colorlinks=true, urlcolor=Blue, linkcolor=Blue, citecolor=Blue}
\usepackage{url}

\usepackage{xspace} %\xspace in \newcommand removes the need for {} to get proper spacing
\usepackage{multicol}
\usepackage{advdate}    % Advancing/saving dates
\usepackage{datetime}   % Dates formatting
\usepackage{datenumber} % Counters for dates
\usepackage{hyphenat} % custom hyphenation ex. $\beta$\hyp{}galactosidase

\usepackage{graphicx}
\usepackage{multicol}
\usepackage{pgfplots,pgfplotstable}
\usepackage{tikz}
\usetikzlibrary{plotmarks}
\pgfplotsset{compat=newest}
\usepackage{filecontents}


%%%biblatex%%%%%%%%%%
\usepackage[style=numeric,natbib=true]{biblatex}
%\usepackage[style=numeric]{biblatex}
\addbibresource{/home/jim/JamesAshe/biblatex_jashe,bib}

\makeatother

\usepackage{babel}
\begin{document}
\newcommand{\me}{Dr.~Ashe\xspace}
\newcommand{\meL}{Dr.~James Ashe\xspace}
\newcommand{\guest}{Dr.~Stephenson\xspace}
\newcommand{\guestL}{Dr.~Ken Stephenson\xspace}
\usdate
\title{Modeling the Thomson Problem with Circle Packings}
\author{Merischia Griffin}
\author{Daniel Villanueva}
\author{James Ashe}
\thanks{Thanks to HBCU-UP of Johnson C.\,Smith for funding this project.}
\thanks{Thanks to Ken Stephenson for serving as an advisor on this project.}
\begin{abstract}
The Thomson problem is longstanding conjecture with applications in
physics and biology. Number 7 on Smale's list of conjectures, the
Thomson problem asks how to find the minimum energy distribution of
$N$ electrons on the surface of the sphere. In general, finding this
minimum is an NP hard problem. Thus progress on this problem must
rely on computational methods. In this project, we use circle packings
to construct a discrete model of the Thomson problem. Circle packings
of minimum configurations are experimentally found to be good approximations
and appear to converge to the actual energy. Furthermore, we demonstrate
how these models can be combinatorial manipulated, and conjecture
how this might be used to compute energy minimums. 
\end{abstract}

\maketitle

\section*{The Thomson Problem}

The Thomson problem asks how to find the minimum energy distribution
of $N$ electrons on the surface of the sphere. The \emph{energy}
of $N$ points, $p_{1},p_{2},\dots,p_{N}$, on the sphere is defined
as \citet{wiki:Thomson}
\[
E={\displaystyle \sum_{i<j}\frac{1}{\left|p_{i}-p_{j}\right|}}\mbox{.}
\]
Number 7 on Smale's list of conjectures, this problem remains one
of the most important unsolved packing problems. Not just important
in mathematics, the Thomson problem is related or equivalent to many
other problems in biology, physics, computer science, chemistry, etc
\citet{wiki:Thomson}. 

Known solutions have been rigorously proven for only a small number
of special cases. In general, finding a solution is an NP hard problem
{[}need reference{]}, and progress has relied on computational methods.
Using such methods, smallest known energy distributions for most points
sets less than 1000 have been found\citet{wales2006structure}. 

\section*{Circle Packing}

Circle packings are configurations of circles with specified patterns
of tangency. The following theorem provides rigidity and existence
of circle packings on the sphere.

\newtheoremstyle{named}{}{}{\itshape}{}{\bfseries}{.}{.5em}{\thmnote{#3 }#1} 
\theoremstyle{named} 
\newtheorem*{namedtheorem}{Theorem} 

\begin{namedtheorem}[Koebe-Andreev-Thurston]

There exists a circle packing on the sphere for any finite triangulation
of the sphere. Furthermore, this circle packing is unique up to M\"{o}bius
transformations.

\end{namedtheorem}

\begin{figure}
\includegraphics[width=6cm]{Thomson_problem_HBCUP_grant_project/images/complex_to_sphere.svg}

\caption{A tangency pattern and its corresponding circle packing.}
\end{figure}

The actual radii for the circle packing guaranteed by the theorem
above can always be found using the \emph{Circle Packing Algorithm}\citet{Col:Ste:MR1975216}.
Moreover we have a feature rich software package, \href{http://www.math.utk.edu/~kens/CirclePack/}{$\tt{CirclePack}$}\footnote{\href{http://www.math.utk.edu/~kens/CirclePack/}{www.math.utk.edu/~kens/CirclePack};
special thanks to Ken Stephenson for developing and maintaining this
freely available software }, for computing and visualizing the lovely results.

Circle packings naturally express a conformal geometry which has allowed
the field of circle packing to be developed into a rich theory of
discrete analytic functions. Many classical complex functions can
be approximated using circles packings, most famously the Riemann
mapping theorem:

\begin{namedtheorem}[Thurston's Circle Packing Conjecture]

The conformal mapping of a bounded simply connected domain onto the
unit disc can be approximated with a circle packing.

\end{namedtheorem}

In summary, we have the following
\begin{enumerate}
\item An essentially unique circle packing for a spherical triangulation
exists.
\item The circle packing can be computationally realized using \href{http://www.math.utk.edu/~kens/CirclePack/}{$\tt{CirclePack}$}.
\item As the number of circles increase, a circle packing analogue will
converge towards its continuous counterpart\citet{Ste:MR2131318}. 
\end{enumerate}

\section*{Our goal}

Repelling points on a flat surface naturally adopt a hexagonal pattern.
However the topology of the sphere precludes a triangulation where
every vertex has 6 neighbors. Try wrapping a sheet of paper around
a basketball --imperfections in your wrapping job are unavoidable.

The Euler characteristic states that there will be 12 such ``imperfections''
where a vertex has 5 neighbors. The truncated icosahedron, a soccer
ball, is a classic example. Thus a combinatorial perspective may shed
new light on this problem. 

We first experimentally demonstrated that circle packings are a good
model of the Thomson problem. Then we demonstrate how manipulating
the combinatorics of circle packings might be used to energy. 

\section*{A discrete Model of the Thomson Problem}

The Thomson problem is well studied, and there is a large catalog
of current minimal energy distributions\citet{Wales:2006:Online}.
These point sets are stored in the interactive web applet, \href{http://thomson.phy.syr.edu/}{\tt{thomsonapplet}}\footnote{\href{http://thomson.phy.syr.edu/}{http://thomson.phy.syr.edu}, \href{mailto:ccecka@seas.harvard.edu}{Cris Cecka},
Primary Author. }, which can also generate random point sets, measure the energy, and
implement a wide array of algorithms. 

To construct a circle packing model for $N$ energy points on the
sphere we did the following: 
\begin{enumerate}
\item Find the Delaunay triangulations for $N$ points using data from \href{http://thomson.phy.syr.edu/}{\tt{thomsonapplet}}. 
\item Compute the circle packing for this triangulation using \href{http://www.math.utk.edu/~kens/CirclePack/}{$\tt{CirclePack}$}.
\item Normalize the packing by applying a M\"{o}bius map such that the
centroid of the circle centers in $\mathbb{R}^{3}$ is at the origin. 
\end{enumerate}
Measure the \emph{error }of this model as, 
\[
error=\frac{\mbox{energy from circle packing-energy}}{\mbox{energy}}
\]
 Figure (XX) shows the error for minimal known distribution for $N$
point. The error from our experiments never exceeded $0.0062$. As
expected the error tended toward $0$ for large values of $N$. It
appears that the error introduced by the discreteness of circle packings
is limited for even small $N$ and inconsequential for larger $N$.

\begin{figure}
\includegraphics{Thomson_problem_HBCUP_grant_project/images/graph_minimal.pdf}

\caption{Error}

\end{figure}


\section*{Combinatorial Manipulaitons}

A Whitehead move, or ``flip'', replaces one interior edge with another
in the union of two edge-sharing faces (see figure XXX). Such a modification
does not change the number of points -only the combinatorics locally.
However these local changes result in a new circle packing and thus
invoke a global geometric change. A common issue with optimization
algorithms is settling into a local but not a global minimum. In the
figure below (figure YYY), a two defects are separated be an area
of degree six vertices. Following a sequence of flips, the defects
are then combined and neutralized.

\begin{figure}
\hfill\includegraphics[height=5cm]{Admin/images/Thomson_410_mesh_thermal.png}\hspace{1cm}\includegraphics[height=5cm]{Admin/images/cpShot-001.jpg}\hfill{}

\caption{\label{fig:An-energy-distribution}An energy distribution modeled
with circle packing.}
\end{figure}

\printbibliography 
\end{document}
